\documentclass[11pt,a4paper]{article}
\usepackage[T1]{fontenc}
\usepackage[utf8]{inputenc}
\usepackage{lmodern}
\usepackage{geometry}
\geometry{margin=1in}
\usepackage{amsmath,amssymb,amsthm,mathtools}
\usepackage{enumitem}
\usepackage{hyperref}
\usepackage{microtype}
\hypersetup{colorlinks=true,linkcolor=blue,urlcolor=blue}

\newcommand{\R}{\mathbb{R}}
\newcommand{\C}{\mathbb{C}}
\newcommand{\N}{\mathbb{N}}
\newcommand{\Z}{\mathbb{Z}}
\newcommand{\E}{\mathbb{E}}
\newcommand{\Prob}{\mathbb{P}}
\newcommand{\one}{\mathbf{1}}
\DeclareMathOperator{\Tr}{Tr}
\DeclareMathOperator{\Sp}{Sp}
\DeclareMathOperator{\diag}{diag}

\title{Ecole Polytechnique--ESPCI 2026\\Mathematics B, MP--MPI\\Concours-Grade Correction}
\author{Detailed English correction}
\date{}

\begin{document}
\maketitle
\tableofcontents
\newpage

\section*{General conventions}
\addcontentsline{toc}{section}{General conventions}
All spectra below are counted with algebraic multiplicity, as in the statement.  For a real symmetric matrix $M\in M_n(\R)$ and a function $f$ defined on the spectrum,
\[
  S_f(M)=\frac1n\sum_{(\lambda,m_\lambda)\in\Sp(M)}m_\lambda f(\lambda).
\]
In Part IV we write $S_n(f)=S_f(X_n)$.  The function $f_k$ is always $x\mapsto x^k$.

\bigskip
This correction gives the details that are usually expected in a concours solution: endpoint cases, multiplicities, convergence estimates, and the probabilistic variance computations are written explicitly.
The appendix records the corresponding Lean formalisation status; it is an
audit aid, not a substitute for the mathematical correction.

\section{Preliminary question}

\paragraph{Question 1.}
The recurrence
\[
  u_0=0,\qquad u_1=1,\qquad u_n=\alpha u_{n-1}-u_{n-2}\quad(n\ge2)
\]
determines the sequence uniquely: if two sequences have the same first two
terms and satisfy the same recurrence, they are equal by induction on $n$.
Thus, in each case below, it is enough to exhibit a sequence with the announced
formula and to check its first two values and the recurrence.

\medskip
\noindent\textbf{Case $|\alpha|>2$.}  Put
\[
  r_+=\frac{\alpha+\sqrt{\alpha^2-4}}2,\qquad
  r_-=\frac{\alpha-\sqrt{\alpha^2-4}}2.
\]
Since $|\alpha|>2$, we have $\alpha^2-4>0$, so $r_+\ne r_-$ and
$r_+-r_-=\sqrt{\alpha^2-4}$.  Moreover $r_+$ and $r_-$ are the two roots of
$r^2-\alpha r+1=0$, hence
\[
  r_\pm^{\,n+2}=\alpha r_\pm^{\,n+1}-r_\pm^{\,n}\qquad(n\ge0).
\]
Therefore the sequence
\[
  v_n=\frac{r_+^n-r_-^n}{r_+-r_-}
\]
satisfies the same recurrence.  Also $v_0=0$ and $v_1=1$.  By uniqueness,
\[
  \boxed{u_n=\frac{r_+^n-r_-^n}{r_+-r_-}
  =\frac{r_+^n-r_-^n}{\sqrt{\alpha^2-4}}\qquad(n\ge0).}
\]

\medskip
\noindent\textbf{Case $|\alpha|<2$.}  There is a unique
$\theta\in(0,\pi)$ such that $\alpha=2\cos\theta$.  Here
$\sin\theta>0$.  Define
\[
  v_n=\frac{\sin(n\theta)}{\sin\theta}.
\]
Then $v_0=0$ and $v_1=1$.  The addition formula gives, for every $n\ge0$,
\[
  \sin((n+2)\theta)
  =2\cos\theta\,\sin((n+1)\theta)-\sin(n\theta).
\]
Dividing by $\sin\theta$ and using $\alpha=2\cos\theta$ shows that
$v_{n+2}=\alpha v_{n+1}-v_n$.  Thus, by uniqueness,
\[
  \boxed{
  \begin{gathered}
    u_n=\dfrac{\sin(n\theta)}{\sin\theta},\\
    \alpha=2\cos\theta,\qquad 0<\theta<\pi.
  \end{gathered}}
\]

\medskip
\noindent\textbf{Case $\alpha=2$.}  Let $v_n=n$.  Then $v_0=0$,
$v_1=1$, and
\[
  2v_{n+1}-v_n=2(n+1)-n=n+2=v_{n+2}.
\]
Hence
\[
  \boxed{u_n=n.}
\]

\medskip
\noindent\textbf{Case $\alpha=-2$.}  Let $v_n=(-1)^{n+1}n$.
Again $v_0=0$ and $v_1=1$.  Moreover
\[
  -2v_{n+1}-v_n
  =-2(-1)^{n+2}(n+1)-(-1)^{n+1}n
  =(-1)^{n+3}(n+2)=v_{n+2}.
\]
Therefore
\[
  \boxed{u_n=(-1)^{n+1}n.}
\]

\section{Part I}

\paragraph{Question 2.}
All estimates below are written for $n\ge1$; changing finitely many initial
terms has no effect on the limit.
Let $f$ be continuous on $[0,1]$.  Since $[0,1]$ is compact, $f$ is uniformly
continuous and bounded.  Set
\[
  M=\sup_{[0,1]}|f|.
\]
We shall use the following elementary Riemann-sum estimate, with the omitted
right endpoint treated explicitly.  Let $m\ge1$, let $h>0$ satisfy
$mh\le1$, and define
\[
  R(m,h)=h\sum_{k=1}^m f(kh).
\]
Then
\[
\begin{aligned}
\left|R(m,h)-\int_0^1 f(t)\,dt\right|
&\le
  \sum_{k=1}^m
  \int_{(k-1)h}^{kh}|f(kh)-f(t)|\,dt
  +\int_{mh}^1 |f(t)|\,dt .
\end{aligned}
\]
For $t\in[(k-1)h,kh]$, both $t$ and $kh$ belong to $[0,1]$ and their distance
is at most $h$.  If
\[
  \omega(\delta)
  =
  \sup\{|f(x)-f(y)|:\ x,y\in[0,1],\ |x-y|\le\delta\},
\]
then $\omega(\delta)\to0$ as $\delta\to0$, and the previous bound gives
\[
  \left|R(m,h)-\int_0^1 f(t)\,dt\right|
  \le mh\,\omega(h)+M(1-mh).
\]
The last term is exactly the contribution of the short interval $[mh,1]$,
whose length tends to $0$ in the applications.  Consequently, whenever
$h_n\to0$ and $m_nh_n\to1$ with $m_nh_n\le1$, one has
\[
  h_n\sum_{k=1}^{m_n} f(kh_n)
  \longrightarrow
  \int_0^1 f(t)\,dt. \tag{2.1}
\]

For the first sequence, take
\[
  h_n=\frac1{n+1},\qquad m_n=n.
\]
Then $h_n\to0$, $m_nh_n=n/(n+1)\to1$, and (2.1) yields
\[
  \frac1{n+1}\sum_{k=1}^n f\left(\frac{k}{n+1}\right)
  \longrightarrow
  \int_0^1 f(t)\,dt.
\]
Since
\[
  v_n
  =
  \frac{n+1}{n}
  \left(
    \frac1{n+1}\sum_{k=1}^n f\left(\frac{k}{n+1}\right)
  \right)
\]
and $(n+1)/n\to1$, we get
\[
  v_n\longrightarrow\int_0^1 f(t)\,dt.
\]

For the second sequence, take
\[
  h_n=\frac2{2n+1},\qquad m_n=n.
\]
Again $h_n\to0$, $m_nh_n=2n/(2n+1)\to1$, and (2.1) gives
\[
  \frac2{2n+1}
  \sum_{k=1}^n
  f\left(\frac{2k}{2n+1}\right)
  \longrightarrow
  \int_0^1 f(t)\,dt.
\]
But
\[
  w_n
  =
  \frac{2n+1}{2n}
  \left(
    \frac2{2n+1}
    \sum_{k=1}^n
    f\left(\frac{2k}{2n+1}\right)
  \right),
\]
and $(2n+1)/(2n)\to1$.  Therefore
\[
  w_n\longrightarrow\int_0^1 f(t)\,dt.
\]
Finally,
\[
  \boxed{\lim_{n\to\infty}v_n=\lim_{n\to\infty}w_n=\int_0^1 f(t)\,dt.}
\]

\paragraph{Question 3a.}
Let $f$ be continuous on $[-2,2]$ and let $M=\sup_{[-2,2]}|f|$.
The only possible singularities of
\[
  \frac{f(x)}{\sqrt{4-x^2}}
\]
are at $x=-2$ and $x=2$.  Near $2$, for $x\in[0,2)$,
\[
  4-x^2=(2-x)(2+x)\ge 2(2-x),
\]
so
\[
  \frac{|f(x)|}{\sqrt{4-x^2}}\le \frac{M}{\sqrt{2}\sqrt{2-x}},
\]
and $1/\sqrt{2-x}$ is integrable near $2$.  The endpoint $-2$ is treated similarly, using
\[
  4-x^2=(2-x)(2+x)\ge 2(x+2)\qquad(x\in(-2,0]).
\]
Thus the improper integral
\[
  I(f)=\frac1\pi\int_{-2}^2\frac{f(x)}{\sqrt{4-x^2}}\,dx
\]
is convergent.

On $[-2+\varepsilon,2-\varepsilon]$ we may use the substitution $x=2\cos\theta$.  Then
\[
  dx=-2\sin\theta\,d\theta,
  \qquad
  \sqrt{4-x^2}=2\sin\theta\quad(0<\theta<\pi).
\]
Letting $\varepsilon\downarrow0$, justified by the convergence just proved, gives
\[
  \boxed{I(f)=\frac1\pi\int_0^\pi f(2\cos\theta)\,d\theta.}
\]

\paragraph{Question 3b.}
For $f_0(x)=1$,
\[
  I(f_0)=\frac1\pi\int_0^\pi1\,d\theta=1.
\]
For $f_1(x)=x$,
\[
  I(f_1)=\frac1\pi\int_0^\pi2\cos\theta\,d\theta=0.
\]
For $f_2(x)=x^2$,
\[
  I(f_2)=\frac1\pi\int_0^\pi4\cos^2\theta\,d\theta
        =\frac4\pi\cdot\frac\pi2=2.
\]
Hence
\[
  \boxed{I(f_0)=1,\qquad I(f_1)=0,\qquad I(f_2)=2.}
\]

\paragraph{Question 3c.}
If $n$ is odd, the function $x^n/\sqrt{4-x^2}$ is odd on $[-2,2]$, so
\[
  \boxed{I(f_n)=0\qquad(n\text{ odd}).}
\]
Let $n=2p$.  Then
\[
  I(f_{2p})=\frac1\pi\int_0^\pi (2\cos\theta)^{2p}\,d\theta.
\]
It remains to compute $\int_0^\pi\cos^{2p}\theta\,d\theta$.  Since $\cos^{2p}$ has period $\pi$,
\[
  \int_0^\pi \cos^{2p}\theta\,d\theta
  =\frac12\int_0^{2\pi}\cos^{2p}\theta\,d\theta.
\]
Using
\[
  \cos\theta=\frac{e^{i\theta}+e^{-i\theta}}2,
\]
we get
\[
  \cos^{2p}\theta
  =2^{-2p}\sum_{j=0}^{2p}\binom{2p}{j}e^{i(2p-2j)\theta}.
\]
The integral over $[0,2\pi]$ kills every nonconstant exponential term.  The constant term corresponds to $j=p$. Hence
\[
  \int_0^{2\pi}\cos^{2p}\theta\,d\theta
  =2\pi\,2^{-2p}\binom{2p}{p},
\]
and therefore
\[
  \int_0^\pi\cos^{2p}\theta\,d\theta
  =\pi\,2^{-2p}\binom{2p}{p}.
\]
Consequently
\[
  I(f_{2p})=\frac{2^{2p}}\pi\cdot \pi\,2^{-2p}\binom{2p}{p}
  =\binom{2p}{p}.
\]
Thus
\[
  \boxed{
  I(f_n)=
  \begin{cases}
  0,& n\text{ odd},\\[1mm]
  \displaystyle \binom{2p}{p},& n=2p.
  \end{cases}}
\]

\paragraph{Question 4a.}
Let $f$ be continuous on $[-2,2]$ and define
\[
  g(t)=f(2\cos(\pi t)),\qquad t\in[0,1].
\]
Then $g$ is continuous on $[0,1]$. Since $U_n$ is uniform on $\{1,\dots,n\}$,
\[
\begin{aligned}
 \E\left(f\left(2\cos\left(\frac{\pi U_n}{n+1}\right)\right)\right)
 &=\frac1n\sum_{k=1}^n f\left(2\cos\left(\frac{\pi k}{n+1}\right)\right)\\
 &=\frac1n\sum_{k=1}^n g\left(\frac{k}{n+1}\right).
\end{aligned}
\]
By Question 2 this converges to
\[
  \int_0^1 g(t)\,dt
  =\frac1\pi\int_0^\pi f(2\cos\theta)\,d\theta
  =\frac1\pi\int_{-2}^2\frac{f(x)}{\sqrt{4-x^2}}\,dx.
\]
Therefore
\[
  \boxed{
  \E\left(f\left(2\cos\left(\frac{\pi U_n}{n+1}\right)\right)\right)
  \longrightarrow I(f).}
\]

\paragraph{Question 4b.}
Fix $y\in[-2,2]$ and define
\[
  \theta_y=\arccos\left(\frac y2\right)\in[0,\pi].
\]
The function $\cos$ is strictly decreasing on $[0,\pi]$.  Thus, for $1\le k\le n$,
\[
  2\cos\left(\frac{\pi k}{n+1}\right)<y
  \iff
  \frac{\pi k}{n+1}>\theta_y.
\]
If $y=-2$, then $\theta_y=\pi$ and the inequality is impossible, so the probability is $0$ for every $n$.  The right-hand integral is also $0$.
If $y=2$, then $\theta_y=0$ and the inequality is true for every $k=1,\dots,n$, so the probability is $1$ for every $n$.  The right-hand integral is also $1$.

Assume now $-2<y<2$.  Put $a_n=(n+1)\theta_y/\pi$.  The number of integers $k\in\{1,\dots,n\}$ with $k>a_n$ is $n-\lfloor a_n\rfloor$.  Hence
\[
  \Prob\left(2\cos\left(\frac{\pi U_n}{n+1}\right)<y\right)
  =\frac{n-\lfloor a_n\rfloor}{n}
  \longrightarrow 1-\frac{\theta_y}{\pi}.
\]
Finally, with $x=2\cos\theta$,
\[
  \frac1\pi\int_{-2}^y\frac{dx}{\sqrt{4-x^2}}
  =\frac1\pi\int_{\theta_y}^{\pi}d\theta
  =1-\frac{\theta_y}{\pi}.
\]
Therefore, for every $y\in[-2,2]$,
\[
  \boxed{
  \Prob\left(2\cos\left(\frac{\pi U_n}{n+1}\right)<y\right)
  \longrightarrow
  \frac1\pi\int_{-2}^y\frac{dx}{\sqrt{4-x^2}}.}
\]

\section{Part II}

\paragraph{Question 5a.}
For $n=2$,
\[
  T_2=\begin{pmatrix}0&1\\1&0\end{pmatrix},
  \qquad
  \chi_2(X)=\det\begin{pmatrix}X&-1\\-1&X\end{pmatrix}=X^2-1.
\]
Thus
\[
  \boxed{\Sp(T_2)=\{-1,1\}.}
\]
For $n=3$,
\[
  XI_3-T_3=\begin{pmatrix}
  X&-1&0\\
  -1&X&-1\\
  0&-1&X
  \end{pmatrix}.
\]
Expanding the determinant gives
\[
  \chi_3(X)=X(X^2-1)-X=X^3-2X=X(X^2-2).
\]
Hence
\[
  \boxed{\Sp(T_3)=\{-\sqrt2,0,\sqrt2\}.}
\]

\paragraph{Question 5b.}
It is convenient to set $\chi_0(X)=1$ and $\chi_1(X)=X$.  For $n\ge2$, expand
\[
  \chi_n(X)=\det(XI_n-T_n)
\]
along the first row.  The first contribution is $X\chi_{n-1}(X)$.  The only other nonzero entry in the first row is $-1$ in column $2$; its cofactor has determinant $-\chi_{n-2}(X)$.  Therefore
\[
  \boxed{\chi_n(X)=X\chi_{n-1}(X)-\chi_{n-2}(X)\qquad(n\ge2).}
\]
In particular this proves the requested recurrence for all $n\ge4$.

\paragraph{Question 5c.}
Let $\alpha\in\C$ with $|\alpha|<2$, and choose a square root $\Delta$ of $4-\alpha^2$.  Since $|\alpha|<2$, one has $\alpha\ne\pm2$, so $\Delta\ne0$.  Put
\[
  r_+=\frac{\alpha+i\Delta}{2},\qquad
  r_-=\frac{\alpha-i\Delta}{2}.
\]
Then
\[
  r_++r_-=\alpha,
  \qquad
  r_+r_-=\frac{\alpha^2+\Delta^2}{4}=1.
\]
Thus $r_+$ and $r_-$ are the two distinct roots of $r^2-\alpha r+1=0$.

For fixed $\alpha$, the numbers $\chi_n(\alpha)$ satisfy
\[
  \chi_n(\alpha)=\alpha\chi_{n-1}(\alpha)-\chi_{n-2}(\alpha),
  \qquad \chi_0(\alpha)=1,
  \qquad \chi_1(\alpha)=\alpha.
\]
The sequence
\[
  q_n=\frac{r_+^{n+1}-r_-^{n+1}}{r_+-r_-}
\]
has the same two initial values and satisfies the same recurrence.  Therefore $\chi_n(\alpha)=q_n$ for all $n\ge0$.  Since $r_+-r_-=i\Delta$, we obtain, for every $n\ge2$,
\[
  \boxed{
  \chi_n(\alpha)=\frac{1}{i\sqrt{4-\alpha^2}}
  \left[
  \left(\frac{\alpha+i\sqrt{4-\alpha^2}}2\right)^{n+1}
  -
  \left(\frac{\alpha-i\sqrt{4-\alpha^2}}2\right)^{n+1}
  \right].}
\]
Changing the sign of the chosen square root exchanges the two terms and changes the sign of the denominator, so the value is independent of this choice.

\paragraph{Question 5d.}
Fix $\alpha\in(-2,2)$ and put $D=\sqrt{4-\alpha^2}$.  Expanding the formula
of Question 5c by the binomial theorem gives
\[
\begin{aligned}
  \chi_n(\alpha)
  &=\frac{(\alpha+iD)^{n+1}-(\alpha-iD)^{n+1}}{iD\,2^{n+1}}\\
  &=\frac1{2^n}\sum_{s=0}^{\lfloor n/2\rfloor}
      (-1)^s\binom{n+1}{2s+1}\alpha^{n-2s}(4-\alpha^2)^s.
\end{aligned}
\]
The last expression defines a polynomial in $\alpha$ and agrees with
$\chi_n(\alpha)$ on the infinite set $(-2,2)$.  Hence
\[
  \chi_n(X)=\frac1{2^n}\sum_{s=0}^{\lfloor n/2\rfloor}
      (-1)^s\binom{n+1}{2s+1}X^{n-2s}(4-X^2)^s
\]
by the polynomial identity principle.
Thus the coefficient of $X^{n-2p}$, for $0\le p\le\lfloor n/2\rfloor$, may be written as
\[
  \frac{(-1)^p4^p}{2^n}
  \sum_{s=p}^{\lfloor n/2\rfloor}
  \binom{n+1}{2s+1}\binom{s}{p}.
\]
This is already an exact expression as a sum of products of binomial coefficients.

It simplifies to the following standard closed form:
\[
  \boxed{\chi_n(X)=\sum_{p=0}^{\lfloor n/2\rfloor}
  (-1)^p\binom{n-p}{p}X^{n-2p}.}
\]
For completeness, we verify the simplification by induction.  The formula is true for $\chi_0=1$ and $\chi_1=X$.  If it holds for $\chi_{n-1}$ and $\chi_{n-2}$, the coefficient of $X^{n-2p}$ in
\[
  X\chi_{n-1}(X)-\chi_{n-2}(X)
\]
is
\[
  (-1)^p\binom{n-1-p}{p}
  -(-1)^{p-1}\binom{n-p-1}{p-1}
  =(-1)^p\binom{n-p}{p},
\]
by Pascal's identity.  The induction follows.

\paragraph{Question 6.}
Let $\theta\in(0,\pi)$ and set $\alpha=2\cos\theta$.  By Question 1 or by Question 5c,
\[
  \chi_n(2\cos\theta)=\frac{\sin((n+1)\theta)}{\sin\theta}.
\]
For
\[
  \theta_k=\frac{k\pi}{n+1},\qquad k=1,\dots,n,
\]
we have $\sin\theta_k\ne0$ and $\sin((n+1)\theta_k)=\sin(k\pi)=0$.  Therefore
\[
  2\cos\left(\frac{k\pi}{n+1}\right)
\]
is a root of $\chi_n$.  These $n$ numbers are distinct because $\cos$ is strictly decreasing on $[0,\pi]$.  Since $\chi_n$ has degree $n$, they are all the roots.  Hence
\[
  \boxed{\Sp(T_n)=\left\{2\cos\left(\frac{k\pi}{n+1}\right):1\le k\le n\right\},}
\]
each eigenvalue having multiplicity $1$.

\paragraph{Question 7.}
By Question 6,
\[
  S_f(T_n)=\frac1n\sum_{k=1}^n
  f\left(2\cos\left(\frac{k\pi}{n+1}\right)\right).
\]
This is exactly the expectation in Question 4a with $U_n$ uniform on $\{1,\dots,n\}$.  Therefore, for every continuous $f$ on $[-2,2]$,
\[
  \boxed{
  \lim_{n\to\infty}S_f(T_n)
  =\frac1\pi\int_{-2}^2\frac{f(x)}{\sqrt{4-x^2}}\,dx.}
\]

\paragraph{Question 8a.}
We have
\[
  T_n(a,b,c)=aI_n+T_n(0,b,c).
\]
If $\mu$ is an eigenvalue of $T_n(0,b,c)$, then $a+\mu$ is an eigenvalue of $T_n(a,b,c)$ with the same algebraic multiplicity.  Equivalently,
\[
  \chi_{T_n(a,b,c)}(X)=\chi_{T_n(0,b,c)}(X-a).
\]
Thus
\[
  \boxed{\Sp(T_n(a,b,c))=a+\Sp(T_n(0,b,c)),}
\]
with multiplicities preserved.

\paragraph{Question 8b.}
Let
\[
  D_n(X)=\det(XI_n-T_n(0,b,c)).
\]
The matrix $XI_n-T_n(0,b,c)$ has diagonal entries $X$, superdiagonal entries $-b$, and subdiagonal entries $-c$.  Expanding as in Question 5b gives
\[
  D_n(X)=X D_{n-1}(X)-bc\,D_{n-2}(X),
  \qquad D_0=1,
  \qquad D_1=X.
\]
The same recurrence and the same initial values are obtained for
\[
  \det(XI_n-T_n(0,bc,1)).
\]
Therefore the two characteristic polynomials are equal.  Together with Question 8a, this gives
\[
  \boxed{\Sp(T_n(a,b,c))=a+\Sp(T_n(0,bc,1)),}
\]
again with algebraic multiplicities.

\paragraph{Question 8c.}
Assume $bc>0$.  Then $b$ and $c$ have the same sign.  Set
\[
  r=\sqrt{\frac cb}>0,
  \qquad
  D=\diag(1,r,r^2,\dots,r^{n-1}).
\]
A direct computation gives
\[
  D^{-1}T_n(0,b,c)D=s\sqrt{bc}\,T_n,
\]
where $s=1$ if $b,c>0$ and $s=-1$ if $b,c<0$.  Indeed the new superdiagonal and subdiagonal entries are respectively $br$ and $c/r$, both equal to $s\sqrt{bc}$.

Thus $T_n(a,b,c)$ is similar to
\[
  aI_n+s\sqrt{bc}\,T_n.
\]
By Question 6 its eigenvalues are
\[
  a+2s\sqrt{bc}\cos\left(\frac{k\pi}{n+1}\right),
  \qquad k=1,
  \dots,n.
\]
If $s=-1$, this list is the same as the list with $s=1$, because replacing $k$ by $n+1-k$ changes the sign of the cosine.  Hence, in all cases $bc>0$,
\[
  \boxed{\lambda_k=a+2\sqrt{bc}\cos\left(\frac{k\pi}{n+1}\right),
  \qquad k=1,\dots,n.}
\]
The eigenvalues are real and simple.

\paragraph{Question 9a.}
For $bc>0$, Question 8c gives
\[
  S_f(T_n(a,b,c))=\frac1n\sum_{k=1}^n
  f\left(a+2\sqrt{bc}\cos\left(\frac{k\pi}{n+1}\right)\right).
\]
Define $g$ on $[-2,2]$ by
\[
  g(x)=f(a+\sqrt{bc}\,x).
\]
If $f$ is continuous on $\R$, then $g$ is continuous on $[-2,2]$.  Applying Question 7 to $g$ gives
\[
\begin{aligned}
 \lim_{n\to\infty}S_f(T_n(a,b,c))
 &=\frac1\pi\int_{-2}^2\frac{g(x)}{\sqrt{4-x^2}}\,dx\\
 &=\frac1\pi\int_{-2}^2\frac{f(a+\sqrt{bc}\,x)}{\sqrt{4-x^2}}\,dx.
\end{aligned}
\]
Thus
\[
  \boxed{\lim_{n\to\infty}S_f(T_n(a,b,c))
  =\frac1\pi\int_{-2}^2\frac{f(a+\sqrt{bc}\,x)}{\sqrt{4-x^2}}\,dx.}
\]

\paragraph{Question 9b.}
Let
\[
  \lambda_{n,k}=a+2\sqrt{bc}\cos\left(\frac{k\pi}{n+1}\right),
  \qquad 1\le k\le n.
\]
The sequence $k\mapsto\lambda_{n,k}$ is strictly decreasing.  Define
\[
  L_-=a-2\sqrt{bc},
  \qquad
  L_+=a+2\sqrt{bc}.
\]

As in the statement, let
\[
  q_n(y)=\#\{1\le k\le n:\ \lambda_{n,k}\le y\},
\]
the number of eigenvalues of $T_n(a,b,c)$ in $(-\infty,y]$, counted with
multiplicity.

If $y<L_-$, then $q_n(y)=0$ for every $n$.  If $y\ge L_+$, then $q_n(y)=n$ for every $n$.

Assume now $L_-<y<L_+$.  Put
\[
  \theta_y=\arccos\left(\frac{y-a}{2\sqrt{bc}}\right)\in(0,\pi).
\]
Then
\[
  \lambda_{n,k}\le y
  \iff
  \frac{k\pi}{n+1}\ge \theta_y.
\]
Hence
\[
  q_n(y)=\#\left\{1\le k\le n:
  k\ge \frac{n+1}{\pi}\theta_y\right\}
  =n-\left\lceil\frac{n+1}{\pi}\theta_y\right\rceil+1,
\]
and dividing by $n$ gives
\[
  \frac{q_n(y)}n\longrightarrow 1-\frac{\theta_y}{\pi}.
\]
Equivalently,
\[
  1-\frac{\theta_y}{\pi}
  =\frac1\pi\int_{-2}^{(y-a)/\sqrt{bc}}
  \frac{dx}{\sqrt{4-x^2}}.
\]
Therefore the clean global statement is
\[
  \boxed{\frac{q_n(y)}n\longrightarrow F(y),}
\]
where
\[
  F(y)=
  \begin{cases}
  0,& y\le a-2\sqrt{bc},\\[1mm]
  \displaystyle \frac1\pi\int_{-2}^{(y-a)/\sqrt{bc}}
  \frac{dx}{\sqrt{4-x^2}},& a-2\sqrt{bc}<y<a+2\sqrt{bc},\\[3mm]
  1,& y\ge a+2\sqrt{bc}.
  \end{cases}
\]
Consequently, whenever $F(y)>0$,
\[
  \boxed{q_n(y)\sim nF(y).}
\]
At the lower edge $y\le a-2\sqrt{bc}$ the count is identically zero, so one should not state a nonzero equivalent there.

\section{Part III: constructive Weierstrass approximation}

The Weierstrass approximation theorem is not used in this part.  It is proved from the polynomials
\[
  Q_n(X)=(1-X^n)^{2^n},
  \qquad
  P_n(X)=Q_n\left(\frac{1-X}{2}\right).
\]

\paragraph{Question 10a.}
Let $0\le\kappa<1/2$.

On $[0,\kappa]$, for $n\ge1$ and $0\le x\le\kappa$,
\[
  0\le x^n\le\kappa^n.
\]
Using $1-(1-t)^m\le mt$ for $0\le t\le1$ and $m\ge1$, we obtain
\[
  0\le 1-Q_n(x)
  =1-(1-x^n)^{2^n}
  \le 2^n x^n
  \le (2\kappa)^n.
\]
Since $2\kappa<1$, this tends uniformly to $0$.  Thus
\[
  Q_n\longrightarrow1
  \quad\text{uniformly on }[0,\kappa].
\]

On $[1-\kappa,1]$, we have $x^n\ge(1-\kappa)^n$.  Hence
\[
  0\le Q_n(x)
  =(1-x^n)^{2^n}
  \le \left(1-(1-\kappa)^n\right)^{2^n}.
\]
Using $1-t\le e^{-t}$,
\[
  Q_n(x)
  \le \exp\left(-2^n(1-\kappa)^n\right)
  =\exp\left(-(2(1-\kappa))^n\right).
\]
Since $2(1-\kappa)>1$, this upper bound tends to $0$.  Thus
\[
  Q_n\longrightarrow0
  \quad\text{uniformly on }[1-\kappa,1].
\]

\paragraph{Question 10b.}
Let $0<\eta\le1$.  For $x\in[-1,1]$ put
\[
  t=\frac{1-x}{2}.
\]
Then $P_n(x)=Q_n(t)$.

If $x\in[\eta,1]$, then
\[
  0\le t\le\frac{1-\eta}{2}=:\kappa<\frac12,
\]
so $Q_n(t)\to1$ uniformly by Question 10a.  Since $H(x)=1$ on $[\eta,1]$, this gives uniform convergence to $H$ there.

If $x\in[-1,-\eta]$, then
\[
  \frac{1+\eta}{2}\le t\le1.
\]
Writing $\kappa=(1-\eta)/2<1/2$, this interval is contained in $[1-\kappa,1]$, so $Q_n(t)\to0$ uniformly by Question 10a.  Since $H(x)=0$ on $[-1,-\eta]$, this gives uniform convergence to $H$ there.

Therefore
\[
  \boxed{P_n\longrightarrow H
  \quad\text{uniformly on }[-1,1]\setminus[-\eta,\eta].}
\]

\paragraph{Question 11.}
Assume first that $f(-1)=0$.  Since $f$ is continuous on the compact interval $[-1,1]$, it is uniformly continuous.  Choose $\delta>0$ such that
\[
  |x-y|\le\delta\quad\Longrightarrow\quad |f(x)-f(y)|\le\varepsilon.
\]
Choose a partition
\[
  -1=t_0<t_1<\cdots<t_N<t_{N+1}=1
\]
with mesh at most $\delta$, and set
\[
  c_i=t_i,
  \qquad
  a_i=f(t_i)-f(t_{i-1})
  \qquad(1\le i\le N).
\]
Then $-1<c_1<\cdots<c_N<1$, and
\[
  |a_i|=|f(t_i)-f(t_{i-1})|\le\varepsilon.
\]
Define
\[
  S(x)=\sum_{i=1}^N a_iH(x-c_i).
\]
If $x\in[t_j,t_{j+1})$ with $0\le j\le N$, then exactly the jumps $c_1,
\dots,c_j$ have occurred, so
\[
  S(x)=\sum_{i=1}^j a_i=f(t_j)-f(t_0)=f(t_j),
\]
because $f(t_0)=f(-1)=0$.  Hence
\[
  |f(x)-S(x)|=|f(x)-f(t_j)|\le\varepsilon.
\]
At $x=1$, one has $S(1)=f(t_N)$ and therefore
\[
  |f(1)-S(1)|=|f(1)-f(t_N)|\le\varepsilon.
\]
Thus
\[
  \boxed{\forall x\in[-1,1],\qquad
  \left|f(x)-\sum_{i=1}^N a_iH(x-c_i)\right|\le\varepsilon,}
\]
with $a_i\in[-\varepsilon,\varepsilon]$.

\paragraph{Question 12.}
We first prove the result for $f(-1)=0$, then remove this assumption.

Let $\varepsilon>0$ and assume $f(-1)=0$.  Apply Question 11 with tolerance $\varepsilon/4$.  We obtain points
\[
  -1<c_1<\cdots<c_N<1
\]
and coefficients $a_i$ with $|a_i|\le\varepsilon/4$ such that
\[
  \left|f(x)-S(x)\right|\le\frac\varepsilon4,
  \qquad
  S(x)=\sum_{i=1}^N a_iH(x-c_i).
\]
For each $i$, choose a positive number $\rho_i$ such that
\[
  \frac{x-c_i}{\rho_i}\in[-1,1]
  \qquad\text{for all }x\in[-1,1].
\]
For example, one may take
\[
  \rho_i=1+|c_i|.
\]
Then
\[
  H\left(\frac{x-c_i}{\rho_i}\right)=H(x-c_i),
\]
because $\rho_i>0$.

Choose $\eta>0$ small enough that the intervals
\[
  [c_i-\rho_i\eta,c_i+\rho_i\eta]
\]
are pairwise disjoint.  This is possible because the points $c_i$ are distinct.  By Question 10b, for each fixed $i$,
\[
  P_m\left(\frac{x-c_i}{\rho_i}\right)
  \longrightarrow
  H\left(\frac{x-c_i}{\rho_i}\right)=H(x-c_i)
\]
uniformly on the set of $x\in[-1,1]$ such that $|x-c_i|\ge \rho_i\eta$.

Let
\[
  A=\sum_{i=1}^N |a_i|.
\]
If $A=0$, then $S=0$ and the proof is immediate.  Otherwise choose $m$ large enough that, simultaneously for all $i$ and all $x\in[-1,1]$ with $|x-c_i|\ge \rho_i\eta$,
\[
  \left|P_m\left(\frac{x-c_i}{\rho_i}\right)-H(x-c_i)\right|
  \le \frac{\varepsilon}{4A}.
\]
Define the polynomial
\[
  R_m(X)=\sum_{i=1}^N a_iP_m\left(\frac{X-c_i}{\rho_i}\right).
\]
Fix $x\in[-1,1]$.  Because the intervals $[c_i-\rho_i\eta,c_i+\rho_i\eta]$ are pairwise disjoint, $x$ belongs to at most one of them.
For all indices for which $|x-c_i|\ge \rho_i\eta$, the total error is bounded by
\[
  \sum_i |a_i|\frac{\varepsilon}{4A}\le\frac\varepsilon4.
\]
For the possible exceptional index $j$, we use only the elementary bound
\[
  0\le P_m(t)\le1\qquad(-1\le t\le1),
\]
which follows from $0\le(1-t)/2\le1$.  Hence
\[
  \left|P_m\left(\frac{x-c_j}{\rho_j}\right)-H(x-c_j)\right|\le1,
\]
and its contribution is at most $|a_j|\le\varepsilon/4$.  Thus
\[
  |R_m(x)-S(x)|\le\frac\varepsilon2.
\]
Together with $|f(x)-S(x)|\le\varepsilon/4$, this gives
\[
  |f(x)-R_m(x)|\le\frac{3\varepsilon}{4}<\varepsilon.
\]
This proves the result when $f(-1)=0$.

For a general continuous $f$, apply the previous case to
\[
  g(x)=f(x)-f(-1),
\]
which satisfies $g(-1)=0$.  There is a polynomial $R$ such that
\[
  |g(x)-R(x)|\le\varepsilon
  \qquad(x\in[-1,1]).
\]
Then
\[
  P(X)=R(X)+f(-1)
\]
satisfies
\[
  \boxed{\forall x\in[-1,1],\qquad |f(x)-P(x)|\le\varepsilon.}
\]

\medskip
\noindent\emph{Scaling note.}  The official hint proposes the unscaled terms
$P_m(X-c_i)$.  Taken literally, Question 10b does not control them on all of
$[-1,1]$, because $x-c_i$ can leave $[-1,1]$.  The rescaling above repairs
this domain gap: $P_m((X-c_i)/\rho_i)$ always receives an argument in
$[-1,1]$.  Since $\rho_i>0$, the target step is unchanged:
$H((x-c_i)/\rho_i)=H(x-c_i)$.  This is one of the places where formal
verification forced a hidden domain condition into the open.

\section{Part IV}

We recall the probabilistic assumptions from the statement.  The
upper-triangular family $(W_{i,j})_{1\le i\le j}$ consists of independent,
identically distributed, integer-valued random variables such that
\[
  \E(W_{1,1})=0,
  \qquad
  \E(W_{1,1}^2)=1,
  \qquad
  \E(|W_{1,1}|^k)<\infty\quad(k\ge0).
\]
The array is extended symmetrically by $W_{j,i}=W_{i,j}$ and
\[
  X_n=\frac1{\sqrt n}(W_{i,j})_{1\le i,j\le n}.
\]

For $k\ge0$, recall
\[
  f_k(x)=x^k,
  \qquad
  \Sigma(f)=\frac1{2\pi}\int_{-2}^2 f(x)\sqrt{4-x^2}\,dx.
\]
For later use, if $c\in\R$ and $j\ge0$, the function
\[
  x\longmapsto c\,x^j\sqrt{4-x^2}
\]
is continuous on the compact interval $[-2,2]$, hence integrable there.  Thus
finite linear combinations of monomials may be integrated term by term against
the semicircle weight.

For every integer $k\ge0$, hypothesis $(H_k)$ is the pair of limits
\[
  \boxed{
  \frac1n\E\!\left[\Tr(X_n^k)\right]\longrightarrow\Sigma(f_k),
  \qquad
  \frac1{n^2}\E\!\left[\Tr(X_n^k)^2\right]\longrightarrow\Sigma(f_k)^2.}
\]
Question 13c proves these limits for $k=0,1,2$; the remainder of the paper
assumes them for every $k$.

\paragraph{Question 13a.}
For every $\omega$, the matrix $X_n(\omega)$ is real symmetric.  Hence it is diagonalizable over $\R$ in an orthonormal basis.  If its eigenvalues are $\lambda_1,\dots,\lambda_n$, repeated with multiplicity, then the eigenvalues of $X_n(\omega)^k$ are $\lambda_1^k,\dots,\lambda_n^k$.  Therefore
\[
  \Tr(X_n(\omega)^k)=\sum_{j=1}^n\lambda_j^k.
\]
On the other hand,
\[
  S_n(f_k)(\omega)=\frac1n\sum_{j=1}^n\lambda_j^k.
\]
Thus, as random variables,
\[
  \boxed{S_n(f_k)=\frac1n\Tr(X_n^k).}
\]

\paragraph{Question 13b.}
If $k$ is odd, the integrand $x^k\sqrt{4-x^2}$ is odd on $[-2,2]$, so
\[
  \boxed{\Sigma(f_k)=0\qquad(k\text{ odd}).}
\]
Let $k=2p$.  With $x=2\cos\theta$,
\[
  dx=-2\sin\theta\,d\theta,
  \qquad
  \sqrt{4-x^2}=2\sin\theta,
\]
and therefore
\[
\begin{aligned}
 \Sigma(f_{2p})
 &=\frac1{2\pi}\int_0^\pi(2\cos\theta)^{2p}\,4\sin^2\theta\,d\theta\\
 &=\frac{2^{2p+1}}\pi\int_0^\pi\cos^{2p}\theta\,\sin^2\theta\,d\theta.
\end{aligned}
\]
From Question 3c,
\[
  A_p:=\int_0^\pi\cos^{2p}\theta\,d\theta
  =\pi\,4^{-p}\binom{2p}{p}.
\]
Since
\[
  \int_0^\pi\cos^{2p}\theta\sin^2\theta\,d\theta
  =A_p-A_{p+1},
\]
and
\[
  \frac{A_{p+1}}{A_p}=\frac{2p+1}{2p+2},
\]
we obtain
\[
  A_p-A_{p+1}=\frac{A_p}{2p+2}.
\]
Thus
\[
\begin{aligned}
 \Sigma(f_{2p})
 &=\frac{2^{2p+1}}\pi\cdot \frac{1}{2p+2}
   \pi\,4^{-p}\binom{2p}{p}\\
 &=\frac1{p+1}\binom{2p}{p}.
\end{aligned}
\]
Consequently
\[
  \boxed{
  \Sigma(f_k)=
  \begin{cases}
  0,& k\text{ odd},\\[1mm]
  \displaystyle \frac1{p+1}\binom{2p}{p},& k=2p.
  \end{cases}}
\]
The even moments are the Catalan numbers.

\paragraph{Question 13c.}
We prove $(H_k)$ for $k=0,1,2$.

\medskip
\noindent\textbf{Case $k=0$.}  Since $X_n^0=I_n$,
\[
  \Tr(X_n^0)=n.
\]
Therefore
\[
  \frac1n\E(\Tr(X_n^0))=1=\Sigma(f_0),
\]
and
\[
  \frac1{n^2}\E(\Tr(X_n^0)^2)=\frac1{n^2}n^2=1=\Sigma(f_0)^2.
\]
Thus $(H_0)$ holds.

\medskip
\noindent\textbf{Case $k=1$.}  We have
\[
  \Tr(X_n)=\frac1{\sqrt n}\sum_{i=1}^n W_{i,i}.
\]
Since the variables have mean zero,
\[
  \frac1n\E(\Tr(X_n))=0=\Sigma(f_1).
\]
Furthermore, by independence and mean zero, the cross terms vanish:
\[
\begin{aligned}
 \E(\Tr(X_n)^2)
 &=\frac1n\E\left(\sum_{i=1}^n W_{i,i}\right)^2\\
 &=\frac1n\sum_{i=1}^n\E(W_{i,i}^2)
 =\frac1n\cdot n=1.
\end{aligned}
\]
Hence
\[
  \frac1{n^2}\E(\Tr(X_n)^2)=\frac1{n^2}\longrightarrow0=\Sigma(f_1)^2.
\]
Thus $(H_1)$ holds.

\medskip
\noindent\textbf{Case $k=2$.}  Since $X_n$ is symmetric,
\[
  \Tr(X_n^2)=\sum_{i,j=1}^n (X_n)_{i,j}^2.
\]
Using the definition of $X_n$,
\[
  \Tr(X_n^2)
  =\frac1n\left(\sum_{i=1}^n W_{i,i}^2
  +2\sum_{1\le i<j\le n}W_{i,j}^2\right).
\]
Thus
\[
\begin{aligned}
 \E(\Tr(X_n^2))
 &=\frac1n\left(n\E(W_{1,1}^2)+2\binom n2\E(W_{1,1}^2)\right)\\
 &=\frac1n\left(n+n(n-1)\right)=n.
\end{aligned}
\]
Therefore
\[
  \frac1n\E(\Tr(X_n^2))=1=\Sigma(f_2).
\]

It remains to prove the second limit.  Put
\[
  Y_{i,j}=W_{i,j}^2,
  \qquad
  A_n=\sum_{i=1}^nY_{i,i}+2\sum_{1\le i<j\le n}Y_{i,j}.
\]
Then
\[
  \Tr(X_n^2)=\frac{A_n}{n},
  \qquad
  \E(A_n)=n^2.
\]
Since $|W_{1,1}|^4$ has finite expectation, $Y_{1,1}$ has finite variance.  Let
\[
  \sigma_Y^2=\operatorname{Var}(Y_{1,1})<\infty.
\]
By independence of the variables $W_{i,j}$ for $i\le j$,
\[
  \operatorname{Var}(A_n)
  =n\sigma_Y^2+4\binom n2\sigma_Y^2
  =O(n^2).
\]
Consequently
\[
  \E(A_n^2)=\E(A_n)^2+\operatorname{Var}(A_n)=n^4+O(n^2).
\]
Hence
\[
  \frac1{n^2}\E(\Tr(X_n^2)^2)
  =\frac1{n^2}\E\left(\frac{A_n^2}{n^2}\right)
  =\frac{\E(A_n^2)}{n^4}
  \longrightarrow1=\Sigma(f_2)^2.
\]
Thus $(H_2)$ holds.

In the remainder of the part, as stated in the problem, we assume $(H_k)$ for every $k\ge0$.

\paragraph{Question 14.}
Let $k\ge0$, $B>0$, and
\[
  g_{k,B}(x)=|x|^k\one_{|x|>B}.
\]
For every real $x$,
\[
  |x|^k\one_{|x|>B}
  \le \frac{|x|^{2k}}{B^k}.
\]
Indeed, on $\{|x|>B\}$ this is equivalent to $B^k\le |x|^k$, and outside this set the left-hand side is zero.  Therefore, for every $\omega$,
\[
  S_n(g_{k,B})(\omega)
  \le \frac1{B^k}S_n(f_{2k})(\omega).
\]
The random variable $S_n(g_{k,B})$ is nonnegative.  Markov's inequality gives
\[
\begin{aligned}
 \Prob(S_n(g_{k,B})\ge\varepsilon)
 &\le \frac1\varepsilon\E(S_n(g_{k,B}))\\
 &\le \frac1{\varepsilon B^k}\E(S_n(f_{2k})).
\end{aligned}
\]
Hence
\[
  \boxed{\Prob(S_n(g_{k,B})\ge\varepsilon)
  \le \frac{\E(S_n(f_{2k}))}{\varepsilon B^k}.}
\]

\paragraph{Question 15.}
The question asks for $B>4$, but a direct comparison reaches the natural
threshold $B>2$.  Fix $k\ge0$, $B>2$, and $\varepsilon>0$.  For every
integer $\ell$ with $2\ell\ge k$, and every $x$ with $|x|>B$,
\[
  |x|^k=|x|^{2\ell}|x|^{k-2\ell}
  \le B^{k-2\ell}|x|^{2\ell}.
\]
Therefore
\[
  g_{k,B}(x)\le B^{k-2\ell}f_{2\ell}(x).
\]
Markov's inequality, Question 13a and $(H_{2\ell})$ give
\[
\begin{aligned}
\limsup_{n\to\infty}\Prob(S_n(g_{k,B})\ge\varepsilon)
&\le \frac{B^{k-2\ell}}{\varepsilon}
   \lim_{n\to\infty}\E(S_n(f_{2\ell}))\\
&=\frac{B^{k-2\ell}}{\varepsilon}\Sigma(f_{2\ell})\\
&\le \frac{B^k}{\varepsilon}\left(\frac4{B^2}\right)^\ell.
\end{aligned}
\]
Here $\Sigma(f_{2\ell})=(\ell+1)^{-1}\binom{2\ell}{\ell}\le4^\ell$.
Letting $\ell\to\infty$ proves
\[
  \boxed{\Prob(S_n(g_{k,B})\ge\varepsilon)\longrightarrow0
  \qquad(B>2).}
\]
This is stronger than the requested $B>4$, and $2$ is the edge of the
limiting semicircle support.

\paragraph{Question 16.}
Let
\[
  Z_n=S_n(f_k)=\frac1n\Tr(X_n^k).
\]
By $(H_k)$,
\[
  \E(Z_n)=\frac1n\E(\Tr(X_n^k))\longrightarrow\Sigma(f_k),
\]
and
\[
  \E(Z_n^2)=\frac1{n^2}\E(\Tr(X_n^k)^2)\longrightarrow\Sigma(f_k)^2.
\]
Therefore
\[
  \operatorname{Var}(Z_n)=\E(Z_n^2)-\E(Z_n)^2\longrightarrow0.
\]
Chebyshev's inequality gives
\[
  \Prob(|Z_n-\E Z_n|\ge\varepsilon)
  \le\frac{\operatorname{Var}(Z_n)}{\varepsilon^2}\longrightarrow0.
\]
Thus
\[
  \boxed{\Prob(|S_n(f_k)-\E(S_n(f_k))|\ge\varepsilon)\longrightarrow0.}
\]

\paragraph{Question 17a.}
Fix $B>4$, let $f$ be continuous on $\R$ and zero outside $[-B,B]$, and fix $\varepsilon>0$.

We shall use explicitly the following elementary properties of the spectral
average.  For each fixed matrix, $S_n$ is a positive linear functional on
functions of the eigenvalue variable:
\[
  S_n(af+bg)=aS_n(f)+bS_n(g),\qquad
  f\le g\Rightarrow S_n(f)\le S_n(g),
\]
and it is normalised by $S_n(1)=1$.  Consequently, if $|h|\le g$, then
\[
  |S_n(h)|\le S_n(g).
\]
These identities follow immediately from the definition
$S_n(h)=n^{-1}\sum_{i=1}^n h(\lambda_i)$, with eigenvalues repeated according
to multiplicity.

Apply Question 12 to the function $t\mapsto f(Bt)$ on $[-1,1]$, then compose the approximating polynomial with $x\mapsto x/B$.  Thus there is a polynomial $P\in\R[X]$ such that
\[
  \sup_{|x|\le B}|f(x)-P(x)|\le \eta,
\]
where $\eta>0$ will be chosen below.  We take
\[
  \eta=\frac\varepsilon{16}.
\]
Since the semicircle density is supported on $[-2,2]\subset[-B,B]$, we have
\[
  |\Sigma(f)-\Sigma(P)|\le \eta\Sigma(1)=\eta.
\]
Indeed, on $[-2,2]$ one has $|f-P|\le\eta$ and
$\sqrt{4-x^2}\ge0$, hence
\[
  |\Sigma(f)-\Sigma(P)|
  \le \frac1{2\pi}\int_{-2}^2 |f(x)-P(x)|\sqrt{4-x^2}\,dx
  \le \eta\Sigma(1).
\]
Here $\Sigma(1)=1$ by Question 13b with $p=0$.  Moreover, because $P$ is a
finite linear combination of monomials, linearity of $S_n$, linearity of
expectation, and the first part of $(H_k)$ give
\[
  \E(S_n(P))\longrightarrow \Sigma(P).
\]
Choose $N$ such that, for all $n\ge N$,
\[
  |\E(S_n(P))-\Sigma(P)|\le\eta.
\]

For every real $x$,
\[
  f(x)-P(x)
  =(f(x)-P(x))\one_{|x|\le B}-P(x)\one_{|x|>B},
\]
because $f(x)=0$ when $|x|>B$.  Therefore, for every $n$ and every $\omega$,
\[
  |S_n(f)(\omega)-S_n(P)(\omega)|
  \le \eta+|S_n(P\one_{|x|>B})(\omega)|.
\]
Indeed, the part $(f-P)\one_{|x|\le B}$ has absolute value at most $\eta$,
so positivity and $S_n(1)=1$ bound its spectral average by $\eta$; the
remaining term is the polynomial tail.
For $n\ge N$, we now write
\[
\begin{aligned}
 |S_n(f)-\Sigma(f)|
 &\le |S_n(P)-\E(S_n(P))|\\
 &\quad +|\E(S_n(P))-\Sigma(P)|
       +|\Sigma(P)-\Sigma(f)|
       +|S_n(f)-S_n(P)|\\
 &\le |S_n(P)-\E(S_n(P))|
       +|S_n(P\one_{|x|>B})|+3\eta.
\end{aligned}
\]
Since $3\eta=3\varepsilon/16<\varepsilon/2$, if both
\[
  |S_n(P\one_{|x|>B})|<\frac\varepsilon4
  \quad\text{and}\quad
  |S_n(P)-\E(S_n(P))|<\frac\varepsilon4,
\]
then
\[
  |S_n(f)-\Sigma(f)|<\varepsilon.
\]
Consequently, for every $n\ge N$,
\[
  \boxed{
  \Prob(|S_n(f)-\Sigma(f)|\ge\varepsilon)
  \le
  \Prob(|S_n(P\one_{|x|>B})|\ge\varepsilon/4)
  +
  \Prob(|S_n(P)-\E(S_n(P))|\ge\varepsilon/4).}
\]

\paragraph{Question 17b.}
We prove that the two probabilities on the right-hand side of Question 17a tend to zero.

First, write
\[
  P(x)=\sum_{j=0}^d c_jx^j.
\]
By Question 16 and the union bound, polynomial fluctuations vanish in probability.  Indeed, if
\[
  C_1=\sum_{j=0}^d |c_j|,
\]
then
\[
  |S_n(P)-\E(S_n(P))|
  \le\sum_{j=0}^d |c_j|\,|S_n(f_j)-\E(S_n(f_j))|.
\]
If $C_1=0$, the quantity is zero.  Otherwise, the event
\[
  |S_n(P)-\E(S_n(P))|\ge\varepsilon/4
\]
is contained in the union of the events
\[
  |S_n(f_j)-\E(S_n(f_j))|\ge\frac{\varepsilon}{4\max(1,C_1)}
  \qquad(0\le j\le d),
\]
each of which has probability tending to zero by Question 16.  Hence
\[
  \Prob(|S_n(P)-\E(S_n(P))|\ge\varepsilon/4)\longrightarrow0.
\]

Second, for $|x|>B>1$ and $0\le j\le d$, we have $|x|^j\le |x|^d$.  Therefore
\[
  |P(x)|\one_{|x|>B}
  \le C_1 |x|^d\one_{|x|>B}=C_1 g_{d,B}(x).
\]
Thus
\[
  |S_n(P\one_{|x|>B})|
  \le S_n(|P|\one_{|x|>B})
  \le C_1 S_n(g_{d,B}).
\]
The first inequality is the estimate $|S_n(h)|\le S_n(|h|)$ applied to
$h=P\one_{|x|>B}$, and the second follows from positivity of $S_n$.
If $C_1=0$, the probability is zero.  Otherwise Question 15 gives
\[
\begin{aligned}
 \Prob(|S_n(P\one_{|x|>B})|\ge\varepsilon/4)
 &\le \Prob\left(S_n(g_{d,B})\ge\frac{\varepsilon}{4C_1}\right)\\
 &\longrightarrow0.
\end{aligned}
\]
Combining this with Question 17a yields
\[
  \boxed{\Prob(|S_n(f)-\Sigma(f)|\ge\varepsilon)\longrightarrow0.}
\]
This is the asserted convergence in probability for every continuous test
function supported in $[-B,B]$.  Since every compact subset of $\R$ lies in
such an interval for some $B>4$, the argument covers every $f\in C_c(\R)$.

It also extends to every bounded continuous function.  Question 15 with
$k=0$ gives
\[
  S_n(\one_{|x|>B})\xrightarrow{\Prob}0\qquad(B>2).
\]
Fix $B>4$ and choose a continuous cutoff $\varphi_B$ equal to $1$ on
$[-B,B]$ and vanishing outside a slightly larger compact interval.  Then
$f\varphi_B\in C_c(\R)$ and
\[
  |S_n(f)-S_n(f\varphi_B)|
  \le \|f\|_\infty S_n(\one_{|x|>B}).
\]
Since the semicircle law is supported on $[-2,2]$, one has
$\Sigma(f\varphi_B)=\Sigma(f)$.  Thus
\[
  \boxed{S_n(f)\xrightarrow{\Prob}\Sigma(f)
  \qquad\text{for every }f\in C_b(\R).}
\]

\section{What the problem is secretly about}

The recurrence polynomials in Part II are the Chebyshev polynomials of the
second kind:
\[
  \chi_n(X)=U_n\!\left(\frac X2\right).
\]
If $q\in\C\setminus\{0,1,-1\}$ and $X=q+q^{-1}$, then
\[
  \chi_n(X)=\frac{q^{n+1}-q^{-(n+1)}}{q-q^{-1}}=[n+1]_q,
\]
the usual quantum integer.  The zeros in Question 6 are its root-of-unity
zeros written on the real line.

The matrix $T_n$ is also the adjacency matrix of the path graph $A_n$.
The normalized trace $n^{-1}\Tr(f(T_n))$ samples the whole cosine grid and
tends to the arcsine law.  By contrast, the corner moment
$\langle e_1,T_n^{2p}e_1\rangle$ counts Dyck paths of length $2p$ and,
once $n$ is large compared with $p$, equals the Catalan number $C_p$, a
moment of the semicircle law.  Thus the arcsine law is the bulk density of
states of long paths, while the semicircle law is the limiting spectral
measure seen from an endpoint.  The same path-graph/Chebyshev structure also
appears in Temperley--Lieb and Jones theory, though none of that language is
needed for the exam.

\appendix
\section{Lean formalisation notes}

This appendix records the state of the Lean formalisation that accompanies the
correction.  It is deliberately separated from the main text: a concours
correction should read as mathematics, whereas the Lean repository records
which statements have been mechanised, which hypotheses remain explicit, and
which proof routes were chosen for formal reasons.

\paragraph{Build status.}
The canonical Lean repository builds with
\[
  \texttt{lake build}.
\]
The source tree \texttt{XENS2026MB} contains no \texttt{sorry},
\texttt{admit}, project \texttt{axiom}, or \texttt{unsafe} declaration.  The
audited endpoints checked during the last pass depend only on Lean's standard
foundational axioms \texttt{propext}, \texttt{Classical.choice}, and
\texttt{Quot.sound}.  This means there is no hidden project-specific axiom in
the audited proof graph.

\paragraph{Ticket status.}
Of the 38 development tickets used for the formalisation, 37 are marked
\emph{done}.  The completed tickets include the recurrence, the Riemann-sum and
arcsine calculations, the constructive Weierstrass proof from the explicit
polynomials $Q_n$ and $P_n$, the path-matrix spectrum, the Toeplitz
factorisation for $bc>0$, the semicircle moments, the tail estimates, the
polynomial concentration reduction, and the compact-support convergence
argument.

The former low-degree model gap is now closed for the concrete Mathlib-usable
Wigner model \texttt{WignerInputConcrete}, with measurable integer entries, all
entry moments, centred upper entries, unit second moments, independent upper
entries, identical distribution, and a uniform fourth-moment bound.

\begin{itemize}[leftmargin=2em]
  \item \textbf{T31}: the low-degree moment hypotheses $(H_1)$ and $(H_2)$ are
  now supplied by
  \[
    \texttt{WignerInputConcrete.lowDegreeMomentAPI}.
  \]
  The last
  nontrivial field was the second degree-two trace moment
  \[
    \mathbb E\!\left[\left(\frac1n\operatorname{Tr}(X_n^2)\right)^2\right]
    \longrightarrow 1,
  \]
  proved as follows.  One writes
  \[
    \frac1n\operatorname{Tr}(X_n^2)
    =
    \frac1{n^2}\sum_{1\le i\le j\le n} c_{ij} W_{ij}^2,
    \qquad c_{ij}\in\{1,2\}.
  \]
  The variables \(W_{ij}^2\) are independent over the upper triangle, and the
  uniform fourth-moment bound gives a uniform bound on their variances.  Hence
  the variance of the weighted upper-triangular sum is \(O(n^2)\), so the
  normalized variance is \(O(n^{-2})\).  Since the expectation is already known
  to be \(1\), the second moment tends to \(1\).
\end{itemize}

The former T34 integrability gap is now closed for concrete Wigner inputs.  In
the namespace \texttt{WignerInputConcrete}, the theorem
\texttt{snMonomialIntegrability} derives the monomial spectral-average
integrability and squared integrability assumptions from finite sums and
products of concrete entry moments.  The abstract concentration theorem remains
conditional for a general \texttt{WignerInput}, as intended, but the concrete
supplier is now compiled.

\paragraph{Question 1.}
The Lean proof follows the same uniqueness principle as the main text: once a
candidate sequence has the correct first two values and satisfies the same
order-two recurrence, it is equal to $u_n$.  For $|\alpha|<2$ the proof uses
the real identity
\[
  \sin((m+2)\theta)
  =
  2\cos\theta\,\sin((m+1)\theta)-\sin(m\theta),
\]
rather than complex roots.  This is the route now displayed in the main text
because it is also the shortest concours-friendly verification.

\paragraph{Question 2.}
The Lean development forced the Riemann-sum proof to keep track of the omitted
endpoint intervals and of the harmless normalising factors
$(n+1)/n$ and $(2n+1)/(2n)$.  The main text has therefore been written with an
explicit modulus-of-continuity estimate and a separate bound for the short
interval $[mh,1]$.

\paragraph{Questions 10--12.}
The no-Weierstrass constraint is preserved in Lean.  The formal proof uses only
the explicit polynomials
\[
  Q_n(X)=(1-X^n)^{2^n},
  \qquad
  P_n(X)=Q_n\!\left(\frac{1-X}{2}\right),
\]
then proves uniform step approximation and polynomial replacement.  The scaling
near a jump $c$ uses
\[
  \frac{x-c}{1+|c|}
\]
so that the argument remains in $[-1,1]$ for $x\in[-1,1]$.  Lean also checks
the bounds $0\le P_n\le1$ and the ``at most one bad jump'' estimate.

\paragraph{Question 17.}
The compact-support convergence proof is mechanised conditionally on the
explicit moment and integrability suppliers stated in the theorem signatures.
The formal proof separates the deterministic approximation inequality, the
polynomial tail bound, the Chebyshev concentration step, and the finite-union
argument.  The remaining issue is not the Question 17 argument itself, but the
probabilistic model needed to derive all suppliers from the entries of
$X_n$.

\paragraph{Current conclusion.}
Lean has not found a false statement in the main concours correction.  It has,
however, made the final probabilistic dependency precise: the prose proof may
assume $(H_k)$ after Question 13 as the problem allows, but a first-principles
formalisation of Question 13c needs a stronger Wigner model than the current
placeholder structure.

\end{document}
